\cleardoublepage
\phantomsection              % for hyperlinks/bookmarks
\chapter*{\Huge\bfseries Dedication}
\thispagestyle{noheader}
\addcontentsline{toc}{chapter}{Dedication}

\vspace{1cm}

\begin{RLtext}
بسم اللّه الرحمان الرحيم \\
و الصلاة و السلام على أشرف المرسلين \\[1em]

الحمد للّه الذي ما تمّ جهدٌ ولا خُتم سعيٌ إلا بفضله، \\
فبجزيل عطائه ارتقيت لهذه المرحلة، \\
و بتوفيق منه سبحانه تكلّل جهدي و اجتهادي لعيش هذه اللحظة! \\
لحظة التخرج من رحلتي التعليمية الحافلة بعديد التجارب، \\
لأبدأ حياة الكفاح و العمل، لتوظيف ما اكتسبته لصالح هذه الأمة. \\[2em]

إلى أهل غزة المستضعفين، \\
في حين يتمحو موضوعي حول عنصر أساسي في التغذية، \\
يعانون اليوم هم من المجاعة. \\
فاللهم انصرهم و حرر فلسطين. \\[2em]

لمن حملتني وهنًا على وهنٍ، \\
و لم تتوان يومًا عن دعمي والتضحية لأجلي.. \\
إلى أمي الغالية، التي طالما رأت سعادتها في سعادتي. \\[2em]

لمن علّمني الصبر و التفاؤل و الإحسان، \\
و شجعني في كل خطوةٍ خطوتها، \\
واِحتفى بكل إنجاز صغير لي وكأنّه شيءٌ عظيم، \\
لأبي الذي لم يرجُ مني سوى نجاحي. \\[2em]

إلى روحِ عمي الطيبة الطاهرة، \\
لازالت ذكراك في قلبي حين رجوت نجاحي و تخرجي! \\[2em]

لأخوتي و أحبتي و سندي: \\
إيمان، هشام، سندس و علي، \\
وجودكم جنبي قوة أعتز بها. \\
أسأل الله أن يديمكم بقربي دائماً و يفتح عليكم بكل ما هو جميل! \\[2em]

لزميلتي التي شاركتني الكفاح في هذا المشروع، \\
من سهرت معي الليالي، و شاركتني المخاوف و الضحكات، \\
إليكِ يا زوليخة، صديقتي الغالية و نعم الصحبة الصالحة. \\[2em]

إلى رفيقاتي اللواتي وقفن بجانبي دائماً: \\
خديجة صديقة الروح، \\
غادة طيبة القلب، \\
سمية جميلة الابتسامة، \\
و سهى ملاكي البريء. \\[2em]

لعائلتي الثانية، عائلتي الترجيحية، \\
التي ساندتني طوال سنين الجامعة المضنية. \\[2em]

و الآن إلى نفسي التي ٱمنت بي، \\
بإذن اللّه سأغدو يومًا كما رجوت، بتوفيقٍ من الله ربي! \\[3em]

\hfill أميرة
\end{RLtext}

\pagebreak
